\section{Husimi convex order for spin representations}
\label{sec:spin}
\label{spin:section}

An occupation measurement removes one fermionic mode.  Its probability is
an affine function of the Majorana covariance matrix, while its conditional
state lives in a smaller spin representation.  Two facts make this reduction
compatible with convex order: one probability mixture represents the
occupation probabilities in every rotated mode, and a nonzero occupation
projection of a pure Gaussian state remains Gaussian.  We prove both facts
for the original Fock-space action before applying the conditional comparison.

\subsection{Fock space, Majorana operators, and coherent rays}

Let
\[
 \mathcal F_n=\Lambda^\bullet\C^n,
 \qquad
 \mathcal F_n^\epsilon
 =\bigoplus_{k\equiv\epsilon\; (\mathrm{mod}\,2)}\Lambda^k\C^n,
 \qquad \epsilon\in\{0,1\}.
\]
The increasing occupation wedges form an orthonormal basis, denoted by
$|s\rangle$, $s\in\{0,1\}^n$.  Creation is exterior multiplication,
$a_j^*\xi=e_j\wedge\xi$, and annihilation $a_j$ is its Hilbert adjoint.
Thus $a_ja_k+a_ka_j=0$ and
$a_ja_k^*+a_k^*a_j=\delta_{jk}\Id$.  Define
\begin{equation}
 c_{2j-1}=a_j+a_j^*,\qquad c_{2j}=i(a_j^*-a_j).
 \label{spin:majorana}
\end{equation}
These Majorana operators are self-adjoint and satisfy
\begin{equation}
 c_\alpha c_\beta+c_\beta c_\alpha
 =2\delta_{\alpha\beta}\Id.
 \label{spin:clifford}
\end{equation}
For $z\in\C^{2n}$ write $c(z)=\sum_\alpha z_\alpha c_\alpha$.
The quadratic form relevant to \eqref{spin:clifford} is the complex
\emph{bilinear} form
$b(z,w)=\sum_\alpha z_\alpha w_\alpha$.

We use the positive real Clifford realization of $\Spin(2n)$: it consists
of products of an even number of real unit Clifford vectors.  Its Fock
representation sends such a product to the corresponding product of the
$c(v)$, and is denoted by $g\mapsto U_g$.  Each factor $c(v)$ is unitary
and reverses parity; hence $U_g$ is unitary and preserves
$\mathcal F_n^\epsilon$.  Conjugation gives
\begin{equation}
 U_g c(v)U_g^*=c(R_gv),\qquad R_g\in\SO(2n).
 \label{spin:orthogonal-action}
\end{equation}
The map onto $\SO(2n)$ is surjective.  Indeed,
$c(w)c(v)c(w)=c(2\langle w,v\rangle w-v)$ for a real unit $w$;
products of two such reflections give the plane rotations generating
$\SO(2n)$.  Allowing odd products realizes the other component of
$O(2n)$.  An implementing word has odd length precisely when the
orthogonal transformation reverses orientation.  These elementary
Clifford facts also describe the parity change of its Fock action.

A \emph{pure Gaussian ray} in $\mathcal F_n^\epsilon$ means a ray in the
$\Spin(2n)$ orbit of an occupation vector of parity $\epsilon$.
All occupation vectors of that parity belong to the same orbit: changing
an even number of occupation bits is effected, up to a phase, by the
corresponding product of Majorana operators.  Choose a unit representative
$\Omega_{n,\epsilon}$ of this orbit.  These are the coherent rays of the
basic half-spin representation; see also \cite{hackl2021} for their
geometric description.

Each parity representation is irreducible.  To see this directly, the
occupation projections
\[
 |s\rangle\langle s|
 =\prod_{j:s_j=1}a_j^*a_j\prod_{j:s_j=0}(\Id-a_j^*a_j)
\]
belong to the complex span of even Majorana words.  Multiplying them by
even words which change the occupied bits gives every matrix unit
$|t\rangle\langle s|$ with $|s|\equiv|t|\pmod2$.  Thus the algebra
spanned by the representation acts as the full matrix algebra on each
sector.  In particular, with $D=2^{n-1}$ and normalized Haar measure $dg$,
\begin{equation}
 \int_{\Spin(2n)}|U_g\Omega_{n,\epsilon}\rangle
       \langle U_g\Omega_{n,\epsilon}|\,dg
 =\frac{\Id_{\mathcal F_n^\epsilon}}D.
 \label{spin:resolution}
\end{equation}
The average commutes with the representation and has trace one, which
proves the formula.  For a density operator $\rho$ on this sector let
\[
 Q_\rho(g)=\langle U_g\Omega_{n,\epsilon},
                    \rho U_g\Omega_{n,\epsilon}\rangle.
\]
Consequently $0\le Q_\rho\le1$ and $\int Q_\rho\,dg=D^{-1}$.
All parity indices below are interpreted modulo two.

\begin{theorem}[Half-spin Husimi convex order]
\label{spin:convex-order}
Let $n\ge1$ and $\epsilon\in\{0,1\}$.  For every density operator
$\rho$ on $\mathcal F_n^\epsilon$, every continuous convex function
$\Phi:[0,1]\to\R$, and every unit pure Gaussian vector
$\Omega\in\mathcal F_n^\epsilon$,
\begin{equation}
 \int_{\Spin(2n)}\Phi(Q_\rho(g))\,dg
 \le
 \int_{\Spin(2n)}\Phi(Q_{|\Omega\rangle\langle\Omega|}(g))\,dg.
 \label{spin:order}
\end{equation}
The right-hand side is independent of $\Omega$.  If $\Phi$ is strictly
convex, equality holds if and only if $\rho$ is a pure Gaussian projector.
\end{theorem}

\subsection{One covariance decomposition for all rotated modes}

The first-mode occupation projection and the Majorana covariance are
\begin{equation}
 P=a_1^*a_1=\frac{\Id+i c_1c_2}{2},\qquad
 (\Gamma_\rho)_{\alpha\beta}
 =\frac i2\Tr\bigl(\rho[c_\alpha,c_\beta]\bigr).
 \label{spin:covariance}
\end{equation}
The matrix $\Gamma_\rho$ is real and skew-symmetric, and
$\Gamma_{U_g\rho U_g^*}=R_g\Gamma_\rho R_g^{\mathsf T}$.
Putting $u=R_ge_1$ and $v=R_ge_2$, we have
\begin{equation}
 a_\rho(g):=\Tr(PU_g^*\rho U_g)
 =\frac{1+u^{\mathsf T}\Gamma_\rho v}{2}.
 \label{spin:occupation}
\end{equation}
Thus every first-mode probability is a readout of the same covariance.

\begin{lemma}[Covariance mixture within one parity orbit]
\label{spin:covariance-mixture}
For any density operator $\rho$ on $\mathcal F_n^\epsilon$ there are
unit pure Gaussian vectors $\chi_s$ of parity $\epsilon$ and numbers
$t_s\ge0$, $\sum_st_s=1$, such that
\begin{equation}
 \Gamma_\rho=\sum_st_s\Gamma_{\chi_s},
 \qquad
 a_\rho(g)=\sum_st_sa_{\chi_s}(g)\quad\hbox{for every }g\in\Spin(2n).
 \label{spin:shared-mixture}
\end{equation}
The probabilities $t_s$ are independent of $g$.
\end{lemma}
\begin{proof}
Choose $R\in\SO(2n)$ such that
\[
 R\Gamma_\rho R^{\mathsf T}
 =\operatorname{diag}(\lambda_1J,\ldots,\lambda_nJ),
 \qquad J=\begin{pmatrix}0&1\\-1&0\end{pmatrix}.
\]
The usual real skew-symmetric normal form gives an orthogonal $R$.
If its determinant is negative, reversing one coordinate changes the
sign of one $\lambda_j$ and makes $R$ special orthogonal.  The signed
parameters allow this adjustment within the given parity sector.
Choose a Spin lift and let $\sigma$ be the correspondingly rotated
density operator.  Set $t_s=\langle s,\sigma s\rangle$.  Positivity and
trace one give a probability distribution supported on
$|s|\equiv\epsilon\pmod2$.

For an occupation vector,
\[
 \Gamma_{|s\rangle}
 =\operatorname{diag}((2s_1-1)J,\ldots,(2s_n-1)J).
\]
The diagonal two-by-two blocks of $\Gamma_\sigma$ are determined by the
occupation expectations $\Tr(\sigma a_j^*a_j)$; its other blocks vanish
by construction.  It follows that
$\Gamma_\sigma=\sum_st_s\Gamma_{|s\rangle}$.
Rotating the occupation vectors back gives the first identity in
\eqref{spin:shared-mixture}.  Equation~\eqref{spin:occupation} then gives
the second identity simultaneously for every $g$.
\end{proof}

This lemma decomposes a covariance matrix.  The many-body operator
$\rho$ itself need not be a mixture of Gaussian projectors.  What enters
the convex comparison is the common set of probabilities in
\eqref{spin:shared-mixture}.

\subsection{Pure spinors and conditional Gaussian closure}

For a nonzero $\chi\in\mathcal F_n$ define its Clifford annihilator by
\begin{equation}
 L_\chi=\{z\in\C^{2n}:c(z)\chi=0\}.
 \label{spin:annihilator}
\end{equation}
Equation~\eqref{spin:clifford} implies that $L_\chi$ is isotropic for
$b$, so $\dim L_\chi\le n$.  A spinor with $\dim L_\chi=n$ is called
\emph{pure}.  This is an algebraic condition on the original Fock vector.
For conditioning at the last mode, we use
$\mathcal F_0^0=\C$, $\mathcal F_0^1=0$ and call the unique nonzero
zero-mode ray Gaussian.

\begin{lemma}[Gaussian orbit and occupation conditioning]
\label{spin:gaussian-closure}
The pure-spinor rays in $\mathcal F_n^\epsilon$ are exactly its pure
Gaussian rays.  If $b(z,z)=0$, $\chi$ is pure, and $c(z)\chi\ne0$,
then $c(z)\chi$ is pure.  In particular, if $P\chi\ne0$, then under
the occupation-space identification
$P\mathcal F_n^\epsilon=|1\rangle\otimes\mathcal F_{n-1}^{\epsilon-1}$,
\begin{equation}
 \frac{P\chi}{\|P\chi\|}=|1\rangle\otimes\phi
 \label{spin:occupied-closure}
\end{equation}
for a unit pure Gaussian vector $\phi$ in the remaining $n-1$ modes.
For the vacant projection $\Id-P$, the same conclusion holds with
first factor $|0\rangle$ and remaining parity $\epsilon$.
\end{lemma}
\begin{proof}
First let $L\subset\C^{2n}$ be maximal isotropic.  It has trivial
intersection with $\overline L$: a nonzero $z$ in that intersection would
satisfy $z,\overline z\in L$ and
$b(z,\overline z)=\sum_\alpha|z_\alpha|^2>0$.
Choose a Hermitian orthonormal basis
$z_j=(p_j+iq_j)/\sqrt2$ of $L$, with $p_j,q_j$ real.
The equations $b(z_j,z_k)=0$ and
$\sum_\alpha\overline{(z_j)_\alpha}(z_k)_\alpha=\delta_{jk}$ show
that $p_1,q_1,\ldots,p_n,q_n$ form a real orthonormal frame.
An orthogonal transformation carries the standard frame to this one.
Implement it by a product $W$ of real unit Majorana operators, as above.
It carries
$\operatorname{Span}\{e_{2j-1}+ie_{2j}:1\le j\le n\}$ to $L$.
The corresponding standard Clifford operators are $2a_j$.

Their joint kernel is precisely the vacuum line: in an occupation
expansion, $a_j\xi=0$ removes every coefficient of a wedge containing
$e_j$, and requiring this for all $j$ leaves only the scalar wedge.
Conjugating by $W$ proves that the joint kernel associated with $L$ is
the single line $\C W|0\rangle$.
Its parity is the length of the implementing word modulo two.  Thus
two such lines of the same parity are related by an even word, hence
by $\Spin(2n)$.  Conversely, the vacuum has an $n$-dimensional
annihilator, and orthogonal conjugation preserves that dimension.
This proves the equality of the pure-spinor and Gaussian ray sets in
each sector, including the odd reference orbit.

Now suppose $c(z)\chi\ne0$ and $b(z,z)=0$.  Then
$z\notin L_\chi=L_\chi^\perp$, so
$K=L_\chi\cap z^\perp$ has dimension $n-1$.
For $w\in K$, the Clifford relation gives
$c(w)c(z)\chi=0$, and $c(z)^2=0$ gives the same conclusion for $w=z$.
Therefore the annihilator of $c(z)\chi$ contains
$K\oplus\C z$, an $n$-dimensional isotropic space.  The resulting
spinor is pure.

Both $a_1$ and $a_1^*$ are nonzero scalar multiples of isotropic
Clifford vectors.  If $P\chi=a_1^*a_1\chi\ne0$, both successive
actions are nonzero, so $P\chi$ is pure.  Write its normalization as
in \eqref{spin:occupied-closure}.  Its annihilator contains the
$a_1^*$ direction.  Isotropy forces every annihilating vector to have
zero $a_1$ coefficient.  After subtracting its $a_1^*$ component,
the remaining vector lies in the last $2n-2$ Majorana directions.
Modulo the $a_1^*$ line, the annihilator therefore gives an
$(n-1)$-dimensional isotropic annihilator of $\phi$.
The first part of the proof identifies $\phi$ as Gaussian.  Removing
one occupied mode changes its parity from $\epsilon$ to $\epsilon-1$.
For a nonzero vacant projection, use $\Id-P=a_1a_1^*$ and reverse
the two null-vector actions.  Its annihilator contains the $a_1$
direction; quotienting by that line proves the same statement for
the remaining modes, now without a parity change.
\end{proof}

\subsection{Haar conditioning and strict equality}

\begin{proof}[Proof of Theorem~\ref{spin:convex-order}]
We induct on $n$, simultaneously for the two parity sectors.  For $n=1$
each sector is one-dimensional and the assertion is immediate.  For
$n\ge2$ we may choose the reference ray in the form
\[
 \Omega_{n,\epsilon}=|1\rangle\otimes\Omega_{n-1,\epsilon-1}.
\]
Let $G_n=\Spin(2n)$ and let $H_n=\Spin(2n-2)$ be the subgroup
generated by the last $2n-2$ Majorana directions.  Its action on the
occupied first-mode subspace is $\Id\otimes U_h$.  Indeed, every
remaining odd Majorana acquires a first-mode parity sign in the tensor
identification, and the signs cancel in each even word defining $H_n$.
Right invariance of Haar measure gives
\begin{equation}
 \int_{G_n}\Phi(Q_\rho(g))\,dg
 =\int_{G_n}\int_{H_n}
  \Phi\!\left(\langle U_gU_h\Omega_{n,\epsilon},
                  \rho U_gU_h\Omega_{n,\epsilon}\rangle\right)\,dh\,dg.
 \label{spin:double-haar}
\end{equation}
These are integrals of continuous bounded functions on compact groups.

For fixed $g$, identify the positive compression
$P U_g^*\rho U_gP$ with $|1\rangle\langle1|\otimes A_g$.
Its trace is $a_\rho(g)$.  If this trace is positive, define
$\rho_g=A_g/a_\rho(g)$ on $\mathcal F_{n-1}^{\epsilon-1}$.
The measurement value in the inner integral is exactly
\begin{equation}
 a_\rho(g)\,
 \langle U_h\Omega_{n-1,\epsilon-1},
         \rho_gU_h\Omega_{n-1,\epsilon-1}\rangle.
 \label{spin:conditioning}
\end{equation}
If $a_\rho(g)=0$, positivity and trace zero give $A_g=0$, and the
inner integral is $\Phi(0)$.  In both cases the outer measure remains
the original Haar probability $dg$.

Let $R$ have the Husimi distribution of a pure Gaussian state in the
smaller sector, and put $H_\Phi(a)=\E\Phi(aR)$.
This function is continuous and convex on $[0,1]$.
The induction hypothesis, applied for each fixed $a=a_\rho(g)$,
and Lemma~\ref{spin:covariance-mixture} give
\begin{align}
 \int_{G_n}\Phi(Q_\rho)\,dg
 &\le\int_{G_n}H_\Phi(a_\rho(g))\,dg \notag\\
 &\le\sum_st_s\int_{G_n}H_\Phi(a_{\chi_s}(g))\,dg \notag\\
 &=\int_{G_n}H_\Phi(a_{\Omega_{n,\epsilon}}(g))\,dg
 =\int_{G_n}\Phi(Q_{|\Omega_{n,\epsilon}\rangle
                         \langle\Omega_{n,\epsilon}|})\,dg.
 \label{spin:chain}
\end{align}
The middle inequality is the common-mixture Jensen step of
Lemma~\ref{ext:common-mixture}, with probabilities fixed for all $g$.
The next equality follows from Haar invariance of the Gaussian orbit.
The last follows from Lemma~\ref{spin:gaussian-closure}: every nonzero
conditional state of the Gaussian input is a lower-dimensional
Gaussian state, so the inner integral has the same law at each $g$.
Zero branches give $\Phi(0)$ on both sides.

Suppose now that $\Phi$ is strictly convex.  By
\eqref{spin:resolution}, $\E R=2^{-(n-2)}>0$, and hence $H_\Phi$
is strictly convex.  Consider two positive-weight endpoints of the
covariance mixture with difference $\Delta\ne0$.  Choose a real unit
$u$ with $\Delta u\ne0$ and set
$v=-\Delta u/\|\Delta u\|$.  Skew-symmetry implies that $u,v$ are
orthonormal and $u^{\mathsf T}\Delta v\ne0$.
For $n\ge2$, $\SO(2n)$ is transitive on ordered orthonormal
two-frames: complete either frame to an oriented orthonormal basis
and map one basis to the other.  Thus some $g$ distinguishes the two
endpoint probabilities in \eqref{spin:occupation}.  Continuity
preserves this difference on a nonempty open set, which has positive
Haar measure.  Jensen's inequality in \eqref{spin:chain} is then
strict.

Equality therefore forces all positive-weight endpoints in the
construction of Lemma~\ref{spin:covariance-mixture} to have the same
covariance.  Distinct occupation strings have distinct diagonal
covariances, so exactly one $t_s$ is nonzero.  For a positive operator,
\[
 |\langle s,\sigma t\rangle|^2
 \le\langle s,\sigma s\rangle\langle t,\sigma t\rangle;
\]
this is Cauchy--Schwarz applied to $\sigma^{1/2}|s\rangle$ and
$\sigma^{1/2}|t\rangle$.  The unique nonzero diagonal entry, together
with trace one, forces $\sigma=|s\rangle\langle s|$.
Undoing the Spin rotation shows that $\rho$ is a pure Gaussian
projector.  Such projectors attain equality by Haar invariance.
\end{proof}

\begin{example}[Occupation conditioning and shared probabilities]
\label{spin:example}
In two modes,
$\chi_\theta=\cos\theta\,|00\rangle+\sin\theta\,|11\rangle$
is Gaussian: it is obtained from the vacuum by
$\exp[\theta(a_1^*a_2^*-a_2a_1)]$, whose generator equals
$(c_1c_3-c_2c_4)/2$.
Its first-mode occupation probability is $\sin^2\theta$;
whenever this is nonzero, the remaining conditional state is the
one-mode occupied ray.

In four modes, consider
$\psi=(|0000\rangle+|1111\rangle)/\sqrt2$.
Its covariance is zero: each mode is occupied with probability
$1/2$, and a quadratic Majorana expression cannot connect its two
occupation components.  Thus its covariance mixture uses the two
fixed weights $1/2,1/2$ on the vacuum and the fully occupied state.
For every rotated mode, \eqref{spin:shared-mixture} gives
$a_\psi(g)=1/2$.  The state is non-Gaussian, since a Gaussian
covariance is an orthogonal conjugate of the vacuum covariance
and therefore squares to $-\Id$.  Its first-mode occupied
conditional state is nevertheless the Gaussian vector $|111\rangle$.
This example separates covariance mixing, conditional states, and
the original many-body state.
\end{example}

\subsection{The basic odd-spin representation}

The odd-dimensional spin representation has the same coherent rays
as a half-spin representation in one larger real dimension.  Equality
of the invariant probability measures on those rays transfers the
distribution comparison without altering its normalization.

\begin{corollary}[Odd-spin Husimi convex order]
\label{spin:odd}
For every $m\ge1$, the basic $2^m$-dimensional representation of
$\Spin(2m+1)$ satisfies the conclusion of
Theorem~\ref{spin:convex-order} on its highest-weight coherent orbit.
For a strictly convex test function, equality holds exactly at the
pure coherent projectors of that orbit.
\end{corollary}
\begin{proof}
Set $N=m+1$.  Fix a real unit vector $e\in\R^{2N}$ and embed
$H=\Spin(2N-1)$ as the lift of its stabilizer in $\SO(2N)$.
The restriction of either half-spin representation to $H$ is the
basic spin representation.  One algebraic verification uses
\[
 \operatorname{Cl}_{2N-1}^0(\C)
 \simeq\operatorname{Cl}_{2N-2}(\C)
 \simeq\operatorname{Mat}_{2^{N-1}}(\C).
\]
The induced unital action on a half-spin space of dimension
$2^{N-1}$ is the standard irreducible matrix-module action.
The Clifford identities follow by using products with one fixed
unit generator for the first isomorphism and the creation and
annihilation matrix units for the second.  Under a maximal torus
of $H$, the weights are
$(\pm\tfrac12,\ldots,\pm\tfrac12)$, each with multiplicity one:
they are the restrictions of the half-spin weights, with the last
sign determined by parity.  In particular, a highest-weight ray
is an occupation Gaussian ray.

For completeness we identify the entire orbit.  The proof of
Lemma~\ref{spin:gaussian-closure} identifies maximal isotropic
annihilators with orthogonal complex structures on $\R^{2N}$.
For example, with $z_j=(p_j+iq_j)/\sqrt2$ one may take
$Jp_j=q_j$ and $Jq_j=-p_j$.  This construction is independent
of the unitary basis of the isotropic space, and the two parity
orbits correspond to its two possible orientation components.

Let $J_1,J_2$ lie in the same component.  An element of
$\SO(2N-1)$ fixing $e$ can first send $J_1e$ to $J_2e$.
After this conjugation the two complex structures agree on
$\operatorname{Span}(e,J_2e)$.  On its orthogonal complement
they induce orthogonal complex structures of the same orientation.
Choose adapted oriented orthonormal bases there; the map between
these bases belongs to $\SO(2N-2)$ and conjugates one restricted
complex structure to the other.  Extending it by the identity on
$\operatorname{Span}(e,J_2e)$ fixes $e$ and completes the conjugation.
Thus $H$ is transitive on the entire half-spin coherent orbit.

The pushforwards of Haar probability from $\Spin(2N)$ and from
$H$ coincide: both are invariant probabilities on this compact
transitive $H$-space.  Uniqueness follows by averaging any
continuous function over $H$; the average is constant on the
orbit, so every invariant probability has the same integral.
The Husimi functions are those of the same rays in the same
Hilbert space.  Theorem~\ref{spin:convex-order}, with $n=N$,
therefore gives the asserted inequality and equality criterion.
\end{proof}
